\section*{Broader impact}

\subsection{Intended use and beneficiaries}

The benchmark is designed for researchers developing reliability monitoring,
fault diagnosis, and early-warning systems for scientific agents. It is an
evaluation artifact, not a deployed safety system. Primary beneficiaries
are the ML reliability community and the developers of scientific-agent
platforms who require measurable, reproducible evidence about failure
dynamics before deploying agents in settings where undetected failure has
scientific or operational consequences.

\subsection{Misuse risks}

The benchmark measures failure prediction under controlled injection. It
does not guarantee detection of all failures under unrestricted deployment,
and it cannot certify that a model with high benchmark scores will be safe
in an open-ended scientific environment. Using benchmark scores as a
standalone safety certification, without the held-out transfer tests and
the negative controls that the artifact requires, is a misuse. The warning
model is deliberately a compact baseline; presenting it as an optimal or
production-ready detector is also a misuse.

\subsection{Privacy and security}

The benchmark does not collect or distribute personal data. Tasks are drawn
from public repositories, synthetic scientific workflows, or access-controlled
datasets with documented permission boundaries. The trace contract records
content hashes and structured metadata rather than raw scientific data within
the public release. Sensitive traces are referenced through archival
identifiers with retrieval procedures that respect data-use agreements.

\subsection{Representational limits}

The workload ladder covers repository repair, ML experimentation, signal
detection, and nonstationary-noise tasks. It does not cover clinical
decision-making, legal reasoning, financial trading, or other high-stakes
domains. Failure dynamics measured in this benchmark should not be
extrapolated to those domains without domain-specific validation. The
benchmark uses two model configurations for the primary study and a third
for transfer; results may not generalize to models with substantially
different architectures or training procedures.

\subsection{Mitigations and residual risks}

The benchmark's governance rules—separation of training reward from
adjudication, concealment of fault truth, frozen splits, grouped
uncertainty, and retention of null and negative results—are structural
mitigations against overclaiming, leakage, and post-hoc interpretation.
Residual risks include the possibility that a model exploits an
unidentified weakness in the adjudicator, that a fault injector introduces
an artifact correlated with the mechanism rather than the mechanism itself
(mitigated by the unseen-injector transfer test), and that pilot-to-
confirmatory changes shift the benchmark population in ways not captured
by the protocol deviations log. The maintenance plan commits to versioned
updates with documented changes and deprecation notices.
